\begin{frame}{역사: 어텐션의 3년}(Vaswani et al., 2017)
  ``어텐션''이라는 용어는 트랜스포머보다 \textbf{3년 앞서 기계번역}
  분야에서 등장했다.
  \small
  \begin{tabular}{@{}llp{7.2cm}@{}}
    \toprule
    연도 & 논문 & 한 줄 \\
    \midrule
    2014 & Bahdanau·Cho·Bengio, \emph{Neural Machine Translation by
      Jointly Learning to Align and Translate} & RNN 번역기의
      ``전체 의미를 은닉 상태 하나에 압축'' 문제를, 각 생성 단계에서
      원문을 다시 참고하는 방식으로 해결 \\
    2015 & Luong et al. & \kb{Query·Key·Value} 라는 이름 체계
      (질문·색인·내용) 정립, dot-product 어텐션 제안 \\
    2017 & Vaswani et al., \emph{Attention Is All You Need} &
      ``이 메커니즘 하나면 RNN 본체가 정말 필요한가?'' — 순환을
      통째로 걷어내고 어텐션 레이어만 쌓아 RNN보다 좋은 성능 \\
    \bottomrule
  \end{tabular}
  \vskip0.8em
  기억해둘 한 줄: \textbf{어텐션의 기원은 RNN의 패치였고,
  트랜스포머는 그 논리의 귀결}이다 — 12.2/12.3의 각 구성요소가
  왜 그 형태인지가 훨씬 빨리 이해된다.
\end{frame}
